\section{Related Work}
\label{sec:related}
Table~\ref{tab:related-comparison} summarizes the selected systems discussed below. Their different tasks, interfaces, and evaluation units allow a qualitative comparison, not a ranking or head-to-head benchmark.

\begin{table*}[!t]
\centering
\caption{Representative related work and the scope of its reported evaluation. The studies address different tasks and are not a common benchmark. OPC UA: Open Platform Communications Unified Architecture.}
\label{tab:related-comparison}
\footnotesize
\setlength{\tabcolsep}{5pt}
\begin{tabularx}{\textwidth}{@{}>{\raggedright\arraybackslash}p{0.16\textwidth}>{\raggedright\arraybackslash}p{0.40\textwidth}X@{}}
\toprule
Study & Main focus & Reported setting \\
\midrule
Biskupovic et al.~\cite{biskupovic2025supervision} & Agent-based industrial service selection and reselection & Service adaptability and performance in an industrial cyber--physical system \\
K\"{o}hle et al.~\cite{koehle2026llmdriven} & Module-specific agents and skill-based manufacturing & Nine simulated production challenges \\
Graul et al.~\cite{graul2026aas} & Grounded proposals, documented authorization, and proposed PLC-side checks & Two proof-of-concept scenarios using a simulated OPC UA server \\
Schall~\cite{schall2026safe} & Process-diagram, supervisory control and data acquisition (SCADA), and procedure-grounded verification & Five simulated scenarios and proposal evaluation \\
Vyas et al.~\cite{vyas2026detection} & Multi-agent fault-recovery planning with a process-plant twin & Simulated mixing and stirred-tank reactor studies \\
This work & Scene-bound signal mappings, shared dispatch, and route-specific governance & Simulator and mock-driver probes, scripted missions, and exploratory live-Agent episodes \\
\bottomrule
\end{tabularx}
\end{table*}

\subsection{Industrial Agents and Plant Integration}
Biskupovic et al. study online service selection and reselection in industrial cyber--physical systems \cite{biskupovic2025supervision}. Xia et al. present a digital-twin-connected automation architecture and control generation \cite{xia2025architecture,xia2025control}; other work generates PLC programs \cite{liu2026agents4plc}. In related control research, human instructions have also been incorporated into model-predictive control \cite{wu2025instructmpc}.

Li et al. use structured instructions, retrieval, and iterative verification to generate and check finite-element scripts for Abaqus; that study evaluates engineering-artifact generation rather than runtime authority over process writes \cite{li2026modsolagent}. Other studies examine large language model (LLM) agents, manufacturing skills, and natural-language machine interaction \cite{koehle2026llmdriven,lim2025llmenabled,zhao2026shopfloor,hofmann2025nlc}. Vyas et al. combine multi-agent fault-recovery planning and plant-knowledge retrieval with Digital Process Plant Twin studies of mixing and continuous stirred-tank reactors (CSTRs) \cite{vyas2026detection}.

These studies address agent organization, generated actions, or task execution. The present study does not introduce an agent policy or control law. It examines the runtime boundary that carries production context and direction-specific signal mappings into an operation and applies the required checks before dispatch.

\subsection{Industrial Interfaces and Semantic Commissioning}
Agent-facing interface work includes Model Context Protocol (MCP) access to manufacturing resources \cite{silva2025mcpmanufacturing}, semantic recipes built from Open Platform Communications Unified Architecture (OPC UA) skills \cite{koehle2025semanticrecipes}, industrial-protocol adapters \cite{xavier2026industriConnect}, and tools for binding manufacturing data to a cognitive digital twin \cite{liu2026semanticcommissioning}. These systems address capability discovery, protocol access, or commissioning. Here, configuration binds known protocol drivers to scene-specific read and write nodes. A named operation resolves those bindings through a runtime path that checks authority and write limits. We do not treat MCP or a semantic node name as an authorization mechanism.

\subsection{Governed Operations and Verification}
Prior work also examines plant-facing verification. Graul et al. propose grounding plans in the Asset Administration Shell, with documented human authorization and PLC-side validation of OPC UA Methods \cite{graul2026aas}. Schall verifies proposals against piping and instrumentation diagrams, SCADA data, and standard operating procedures; Lin and Chen describe a policy-governed sensor--manufacturing execution system (MES)--enterprise resource planning (ERP) digital thread \cite{schall2026safe,lin2026agentic}. PLCBench evaluates autonomous agents against commercial PLCs with hardware-in-the-loop process simulation \cite{zhou2026plcbench}. These studies complement the runtime examined here and illustrate established concerns about authorization and safety controls. Our implementation examines how configured signal bindings, route-specific checks, shared dispatch, and separate command and process records operate together at the write boundary. This is a system-integration contribution; the studies use different settings and do not establish comparative safety.

\subsection{Production Context and Outcome Evidence}
MCP standardizes tool discovery and invocation \cite{mcp2026}. For industrial operations, application-layer services must still enforce operation-specific authority and process constraints. ISA-88 and ISA-95 define recipe and production-operation concepts \cite{isa88_2010,isa95_2025}. Run-to-run control relates process observations to subsequent decisions, while human--automation research examines allocation of responsibilities between people and automation \cite{moyne2001r2r,parasuraman2000model}. These studies establish important pieces of industrial-agent integration, but they do not make the boundary invariants used here explicit as one operation contract. In particular, the literature does not jointly expose context as non-authorizing information, direction-specific acquisition/control bindings, post-approval route revalidation, typed separation of command/readback/process evidence, and convergence of task-mediated and direct calls on the same guarded boundary. The present work therefore makes those invariants explicit and evaluates selected application-level admission paths in simulators and mock-driver fixtures. This is a system-integration contribution, not a claim of comparative control efficacy, functional safety, or complete agent-to-device provenance.
